\section{LLM Agent 的评估与实验室}

\subsection{定量评估指标}

企业部署 Agent 之前需要构建一套定量值指标：

\begin{itemize}
    \item \textbf{Task Success Rate}：目标完成率（例如自动整理报告并获得审批）
    \item \textbf{Force Stop Rate}：用户主动中止对话的比例
    \item \textbf{Tool Error Rate}：API/工具出现错误导致 workflow failure
    \item \textbf{Response Divergence}：不同模型版本输出的一致性
\end{itemize}

这些指标通常按 SLA 分层：Critical Task 用 99.9\% success，Exploratory Task 允许更低门槛。

\subsection{Human-in-the-loop 实验室}

Burak 提到企业会设立「Agent Lab」：小规模用户面向真实场景的 testing。流程如下：

\begin{enumerate}
    \item 模拟真实任务，让 agent 在尽可能真实的系统中执行
    \item 安排人工 reviewer 定期 audit output，打标签
    \item 把 label 结果做成 dashboard，触发 retrain / prompt tune
\end{enumerate}

\begin{knowledgebox}{Agent Lab 的组织价值}
Agent Lab 是 bridge between research and product：研究团队可以看到真实 failure，产品团队可以体验新版本，治理团队可以验证安全 guardrails 是否有效。
\end{knowledgebox}

\subsection{本章小结}

评估体系由定量指标与人类实验室组成，保证 agent 既能高效，也能在失败时迅速定位 root cause。

